\section{结论与展望}

在本软件开发文档中，我们从需求分析出发，设计并规划了一个面向全市场、跨资产、支持自校准的金融市场仿真平台。通过系统性的架构设计与模块划分，本项目致力于在以下几方面取得突破：

\begin{itemize}
  \item \textbf{跨资产与跨日能力：}针对现有模拟器主要面向单资产日内交易的局限，我们提出了支持多市场、多资产和集合竞价/连续竞价切换的仿真架构，能够更真实地反映跨资产联动和长周期交易行为。
  \item \textbf{高性能核心内核：}借鉴 ABIDES 消息驱动、离散事件模拟的思想，我们设计了可并行化的事件调度器和高效订单簿数据结构，为大规模仿真提供基础保障。
  \item \textbf{动态自校准机制：}将市场统计特征监测与代理参数调整嵌入仿真循环，实现“在线”校准，减少迭代开销并提升与真实市场的拟合度。
  \item \textbf{丰富的代理框架与交互接口：}提供多种内置代理和自定义接口，并通过程序化 API、命令行、可视化 UI 等手段方便用户测试策略与政策。
  \item \textbf{工程化规范与可扩展性：}采用模块化、分层化设计，结合清晰的接口规范和开发计划，方便团队协作与后续维护扩展。
\end{itemize}

虽然目前开发仍在进行中，但文档已搭建了完整的框架，包括需求说明、系统架构、模块设计、开发计划及开发日志。后续工作将围绕以下方向深入展开：

\begin{itemize}
  \item \textbf{核心实现与优化：}持续迭代核心内核、订单簿和撮合引擎，实现跨日仿真功能并在大规模场景下优化性能。
  \item \textbf{算法与模型扩展：}探索更先进的自校准算法（如强化学习驱动的市场生成模型），引入更多类型的交易代理以模拟复杂行为。
  \item \textbf{交互与可视化完善：}开发完整的 Web UI 或桌面应用，提供实时监控、策略调试和反事实分析工具，提升用户体验。
  \item \textbf{插件化与生态构建：}设计插件系统，使外部研究者能够方便地集成自定义市场规则、代理模型或分析模块，形成开放的社区生态。
  \item \textbf{应用与验证：}通过真实市场数据回放、假设情景测试和策略竞赛等方式验证仿真器的有效性，并与 ABIDES 等现有平台进行对比实验。
\end{itemize}

随着开发的推进，本文档将不断更新，补充具体实现细节、测试结果与案例分析，为最终的软件发布和学术论文撰写提供全面的支撑。我们期望在不久的将来推出一个功能完整、高保真度且易用的全尺寸市场仿真平台，为金融研究、策略开发和政策评估提供强大的工具。